\section{AHSME 1959}
        \begin{enumerate}
        	\item (\textit{AHSME 1959}) Each edge of a cube is increased by $50$\%. The percent of increase of the surface area of the cube is:
$\textbf{(A)}\ 50 \qquad\textbf{(B)}\ 125\qquad\textbf{(C)}\ 150\qquad\textbf{(D)}\ 300\qquad\textbf{(E)}\ 750$


	\item (\textit{AHSME 1959}) Through a point $P$ inside the $\triangle ABC$ a line is drawn parallel to the base $AB$, dividing the triangle into two equal areas. 
If the altitude to $AB$ has a length of $1$, then the distance from $P$ to $AB$ is:
$\textbf{(A)}\ \frac12 \qquad\textbf{(B)}\ \frac14\qquad\textbf{(C)}\ 2-\sqrt2\qquad\textbf{(D)}\ \frac{2-\sqrt2}{2}\qquad\textbf{(E)}\ \frac{2+\sqrt2}{8}$


	\item (\textit{AHSME 1959}) If the diagonals of a quadrilateral are perpendicular to each other, the figure would always be included under the general classification:
$\textbf{(A)}\ \text{rhombus} \qquad\textbf{(B)}\ \text{rectangles} \qquad\textbf{(C)}\ \text{square} \qquad\textbf{(D)}\ \text{isosceles trapezoid}\qquad\textbf{(E)}\ \text{none of these}$


	\item (\textit{AHSME 1959}) If $78$ is divided into three parts which are proportional to $1, \frac13, \frac16,$ the middle part is:
$\textbf{(A)}\ 9\frac13 \qquad\textbf{(B)}\ 13\qquad\textbf{(C)}\ 17\frac13 \qquad\textbf{(D)}\ 18\frac13\qquad\textbf{(E)}\ 26$


	\item (\textit{AHSME 1959}) The value of $\left(256\right)^{.16}\left(256\right)^{.09}$ is:


 $\textbf{(A)}\ 4 \qquad \\ \textbf{(B)}\ 16\qquad \\ \textbf{(C)}\ 64\qquad \\ \textbf{(D)}\ 256.25\qquad \\ \textbf{(E)}\ -16$


	\item (\textit{AHSME 1959}) Given the true statement: If a quadrilateral is a square, then it is a rectangle. 
It follows that, of the converse and the inverse of this true statement is:


 $\textbf{(A)}\ \text{only the converse is true} \qquad \\ \textbf{(B)}\ \text{only the inverse is true }\qquad  \\ \textbf{(C)}\ \text{both are true} \qquad \\ \textbf{(D)}\ \text{neither is true} \qquad \\ \textbf{(E)}\ \text{the inverse is true, but the converse is sometimes true}$


	\item (\textit{AHSME 1959}) The sides of a right triangle are $a$, $a+d$, and $a+2d$, with $a$ and $d$ both positive. The ratio of $a$ to $d$ is:
$\textbf{(A)}\ 1:3 \qquad\textbf{(B)}\ 1:4 \qquad\textbf{(C)}\ 2:1\qquad\textbf{(D)}\ 3:1\qquad\textbf{(E)}\ 3:4$


	\item (\textit{AHSME 1959}) The value of $x^2-6x+13$ can never be less than:


 $\textbf{(A)}\ 4 \qquad \textbf{(B)}\ 4.5 \qquad \textbf{(C)}\ 5\qquad \textbf{(D)}\ 7\qquad \textbf{(E)}\ 13$


	\item (\textit{AHSME 1959}) A farmer divides his herd of $n$ cows among his four sons so that one son gets one-half the herd, 
a second son, one-fourth, a third son, one-fifth, and the fourth son, $7$ cows. Then $n$ is:


 $\textbf{(A)}\ 80 \qquad\textbf{(B)}\ 100\qquad\textbf{(C)}\ 140\qquad\textbf{(D)}\ 180\qquad\textbf{(E)}\ 240$


	\item (\textit{AHSME 1959}) In $\triangle ABC$ with $\overline{AB}=\overline{AC}=3.6$, a point $D$ is taken on $AB$ at a distance $1.2$ from $A$. 
Point $D$ is joined to $E$ in the prolongation of $AC$ so that $\triangle AED$ is equal in area to $ABC$. Then $\overline{AE}$ is:
$\textbf{(A)}\ 4.8 \qquad\textbf{(B)}\ 5.4\qquad\textbf{(C)}\ 7.2\qquad\textbf{(D)}\ 10.8\qquad\textbf{(E)}\ 12.6$


	\item (\textit{AHSME 1959}) The logarithm of $.0625$ to the base $2$ is:
$\textbf{(A)}\ .025 \qquad\textbf{(B)}\ .25\qquad\textbf{(C)}\ 5\qquad\textbf{(D)}\ -4\qquad\textbf{(E)}\ -2$


	\item (\textit{AHSME 1959}) By adding the same constant to $20,50,100$ a geometric progression results. The common ratio is:
$\textbf{(A)}\ \frac53 \qquad\textbf{(B)}\ \frac43\qquad\textbf{(C)}\ \frac32\qquad\textbf{(D)}\ \frac12\qquad\textbf{(E)}\ \frac{1}3$


	\item (\textit{AHSME 1959}) The arithmetic mean (average) of a set of $50$ numbers is $38$. If two numbers, namely, $45$ and $55$, are discarded, the mean of the remaining set of numbers is:
$\textbf{(A)}\ 36.5 \qquad\textbf{(B)}\ 37\qquad\textbf{(C)}\ 37.2\qquad\textbf{(D)}\ 37.5\qquad\textbf{(E)}\ 37.52$


	\item (\textit{AHSME 1959}) Given the set $S$ whose elements are zero and the even integers, positive and negative. 
Of the five operations applied to any pair of elements: (1) addition (2) subtraction 
(3) multiplication (4) division (5) finding the arithmetic mean (average), those elements that only yield elements of $S$ are:
$\textbf{(A)}\ \text{all} \qquad\textbf{(B)}\ 1,2,3,4\qquad\textbf{(C)}\ 1,2,3,5\qquad\textbf{(D)}\ 1,2,3\qquad\textbf{(E)}\ 1,3,5$


	\item (\textit{AHSME 1959}) In a right triangle the square of the hypotenuse is equal to twice the product of the legs. One of the acute angles of the triangle is:
$\textbf{(A)}\ 15^{\circ} \qquad\textbf{(B)}\ 30^{\circ} \qquad\textbf{(C)}\ 45^{\circ} \qquad\textbf{(D)}\ 60^{\circ}\qquad\textbf{(E)}\ 75^{\circ}$


	\item (\textit{AHSME 1959}) The expression$\frac{x^2-3x+2}{x^2-5x+6}\div \frac{x^2-5x+4}{x^2-7x+12},$ when simplified is:
$\textbf{(A)}\ \frac{(x-1)(x-6)}{(x-3)(x-4)} \qquad\textbf{(B)}\ \frac{x+3}{x-3}\qquad\textbf{(C)}\ \frac{x+1}{x-1}\qquad\textbf{(D)}\ 1\qquad\textbf{(E)}\ 2$


	\item (\textit{AHSME 1959}) If $y=a+\frac{b}{x}$, where $a$ and $b$ are constants, and if $y=1$ when $x=-1$, and $y=5$ when $x=-5$, then $a+b$ equals:
$\textbf{(A)}\ -1 \qquad\textbf{(B)}\ 0\qquad\textbf{(C)}\ 1\qquad\textbf{(D)}\ 10\qquad\textbf{(E)}\ 11$


	\item (\textit{AHSME 1959}) The arithmetic mean (average) of the first $n$ positive integers is:
$\textbf{(A)}\ \frac{n}{2} \qquad\textbf{(B)}\ \frac{n^2}{2}\qquad\textbf{(C)}\ n\qquad\textbf{(D)}\ \frac{n-1}{2}\qquad\textbf{(E)}\ \frac{n+1}{2}$


	\item (\textit{AHSME 1959}) With the use of three different weights, namely $1$ lb., $3$ lb., and $9$ lb., how many objects of different weights can be weighed, 
if the objects is to be weighed and the given weights may be placed in either pan of the scale?
$\textbf{(A)}\ 15 \qquad\textbf{(B)}\ 13\qquad\textbf{(C)}\ 11\qquad\textbf{(D)}\ 9\qquad\textbf{(E)}\ 7$


	\item (\textit{AHSME 1959}) It is given that $x$ varies directly as $y$ and inversely as the square of $z$, and that $x=10$ when $y=4$ and $z=14$. Then, when $y=16$ and $z=7$, $x$ equals:
$\textbf{(A)}\ 180\qquad \textbf{(B)}\ 160\qquad \textbf{(C)}\ 154\qquad \textbf{(D)}\ 140\qquad \textbf{(E)}\ 120$


	\item (\textit{AHSME 1959}) If $p$ is the perimeter of an equilateral $\triangle$ inscribed in a circle, the area of the circle is:
$\textbf{(A)}\ \frac{\pi p^2}{3} \qquad\textbf{(B)}\ \frac{\pi p^2}{9}\qquad\textbf{(C)}\ \frac{\pi p^2}{27}\qquad\textbf{(D)}\ \frac{\pi p^2}{81}\qquad\textbf{(E)}\ \frac{\pi p^2\sqrt3}{27}$


	\item (\textit{AHSME 1959}) The line joining the midpoints of the diagonals of a trapezoid has length $3$. If the longer base is $97,$ then the shorter base is:
$\textbf{(A)}\ 94 \qquad\textbf{(B)}\ 92\qquad\textbf{(C)}\ 91\qquad\textbf{(D)}\ 90\qquad\textbf{(E)}\ 89$


	\item (\textit{AHSME 1959}) The set of solutions of the equation $\log_{10}\left( a^2-15a\right)=2$ consists of
$\textbf{(A)}\ \text{two integers } \qquad\textbf{(B)}\ \text{one integer and one fraction}\qquad \textbf{(C)}\ \text{two irrational numbers }\qquad\textbf{(D)}\ \text{two non-real numbers} \qquad\textbf{(E)}\ \text{no numbers, that is, the empty set}$


	\item (\textit{AHSME 1959}) A chemist has m ounces of salt water that is $m$\% salt. How many ounces of salt must he add to make a solution that is $2m$\% salt?


 $\textbf{(A)}\ \frac{m}{100+m} \qquad\textbf{(B)}\ \frac{2m}{100-2m}\qquad\textbf{(C)}\ \frac{m^2}{100-2m}\qquad\textbf{(D)}\ \frac{m^2}{100+2m}\qquad\textbf{(E)}\ \frac{2m}{100+2m}$


	\item (\textit{AHSME 1959}) The symbol $|a|$ means $+a$ if $a$ is greater than or equal to zero, and $-a$ if a is less than or equal to zero; the symbol $<$ means "less than"; 
the symbol $>$ means "greater than."
The set of values $x$ satisfying the inequality $|3-x|<4$ consists of all $x$ such that:
$\textbf{(A)}\ x^2<49 \qquad\textbf{(B)}\ x^2>1 \qquad\textbf{(C)}\ 1<x^2<49\qquad\textbf{(D)}\ -1<x<7\qquad\textbf{(E)}\ -7<x<1$


	\item (\textit{AHSME 1959}) The base of an isosceles triangle is $\sqrt 2$. The medians to the legs intersect each other at right angles. The area of the triangle is:
$\textbf{(A)}\ 1.5 \qquad\textbf{(B)}\ 2\qquad\textbf{(C)}\ 2.5\qquad\textbf{(D)}\ 3.5\qquad\textbf{(E)}\ 4$


	\item (\textit{AHSME 1959}) Which one of the following is  not  true for the equation $ix^2-x+2i=0$, where $i=\sqrt{-1}$
$\textbf{(A)}\ \text{The sum of the roots is 2} \qquad \\ \textbf{(B)}\ \text{The discriminant is 9}\qquad \\ \textbf{(C)}\ \text{The roots are imaginary}\qquad \\ \textbf{(D)}\ \text{The roots can be found using the quadratic formula}\qquad \\ \textbf{(E)}\ \text{The roots can be found by factoring, using imaginary numbers}$


	\item (\textit{AHSME 1959}) In triangle $ABC$, $AL$ bisects angle $A$, and $CM$ bisects angle $C$. Points $L$ and $M$ are on $BC$ and $AB$, respectively. The sides of $\triangle ABC$ are $a$, $b$, and $c$. Then $\frac{AM}{MB}=k\frac{CL}{LB}$ where $k$ is:
$\textbf{(A)}\ 1 \qquad\textbf{(B)}\ \frac{bc}{a^2}\qquad\textbf{(C)}\ \frac{a^2}{bc}\qquad\textbf{(D)}\ \frac{c}{b}\qquad\textbf{(E)}\ \frac{c}{a}$


	\item (\textit{AHSME 1959}) On a examination of $n$ questions a student answers correctly $15$ of the first $20$. Of the remaining questions he answers one third correctly. 
All the questions have the same credit. If the student's mark is 50\%, how many different values of $n$ can there be?
$\textbf{(A)}\ 4 \qquad\textbf{(B)}\ 3\qquad\textbf{(C)}\ 2\qquad\textbf{(D)}\ 1\qquad\textbf{(E)}\ \text{the problem cannot be solved}$


	\item (\textit{AHSME 1959}) $A$ can run around a circular track in $40$ seconds. $B$, running in the opposite direction, meets $A$ every $15$ seconds. 
What is $B$'s time to run around the track, expressed in seconds?
$\textbf{(A)}\ 12\frac12 \qquad\textbf{(B)}\ 24\qquad\textbf{(C)}\ 25\qquad\textbf{(D)}\ 27\frac12\qquad\textbf{(E)}\ 55$


	\item (\textit{AHSME 1959}) A square, with an area of $40$, is inscribed in a semicircle. The area of a square that could be inscribed in the entire circle with the same radius, is:
$\textbf{(A)}\ 80 \qquad\textbf{(B)}\ 100\qquad\textbf{(C)}\ 120\qquad\textbf{(D)}\ 160\qquad\textbf{(E)}\ 200$


	\item (\textit{AHSME 1959}) The length $l$ of a tangent, drawn from a point $A$ to a circle, is $\frac43$ of the radius $r$. The (shortest) distance from A to the circle is:
$\textbf{(A)}\ \frac{1}{2}r \qquad\textbf{(B)}\ r\qquad\textbf{(C)}\ \frac{1}{2}l\qquad\textbf{(D)}\ \frac23l \qquad\textbf{(E)}\ \text{a value between r and l.}$


	\item (\textit{AHSME 1959}) A harmonic progression is a sequence of numbers such that their reciprocals are in arithmetic progression.
Let $S_n$ represent the sum of the first $n$ terms of the harmonic progression; for example $S_3$ represents the sum of 
the first three terms. If the first three terms of a harmonic progression are $3,4,6$, then:
$\textbf{(A)}\ S_4=20 \qquad\textbf{(B)}\ S_4=25\qquad\textbf{(C)}\ S_5=49\qquad\textbf{(D)}\ S_6=49\qquad\textbf{(E)}\ S_2=\frac{1}2 S_4$


	\item (\textit{AHSME 1959}) Let the roots of $x^2-3x+1=0$ be $r$ and $s$. Then the expression $r^2 + s^2$ is:
$\textbf{(A)}\ \text{a positive integer} \qquad\textbf{(B)}\ \text{a positive fraction greater than 1}\qquad\textbf{(C)}\ \text{a positive fraction less than 1}\qquad\textbf{(D)}\ \text{an irrational number}\qquad\textbf{(E)}\ \text{an imaginary number}$


	\item (\textit{AHSME 1959}) The symbol $\ge$ means "greater than or equal to"; the symbol $\le$ means "less than or equal to".
In the equation $(x-m)^2-(x-n)^2=(m-n)^2; m$ is a fixed positive number, and $n$ is a fixed negative number. The set of values x satisfying the equation is:
$\textbf{(A)}\ x\ge 0 \qquad\textbf{(B)}\ x\le n\qquad\textbf{(C)}\ x=0\qquad\textbf{(D)}\ \text{the set of all real numbers}\qquad\textbf{(E)}\ \text{none of these}$


	\item (\textit{AHSME 1959}) The base of a triangle is $80$, and one side of the base angle is $60^\circ$. The sum of the lengths of the other two sides is $90$. The shortest side is:
$\textbf{(A)}\ 45 \qquad\textbf{(B)}\ 40\qquad\textbf{(C)}\ 36\qquad\textbf{(D)}\ 17\qquad\textbf{(E)}\ 12$


	\item (\textit{AHSME 1959}) When simplified the product $\left(1-\frac13\right)\left(1-\frac14\right)\left(1-\frac15\right)\cdots\left(1-\frac1n\right)$ becomes:
$\textbf{(A)}\ \frac1n \qquad\textbf{(B)}\ \frac2n\qquad\textbf{(C)}\ \frac{2(n-1)}{n}\qquad\textbf{(D)}\ \frac{2}{n(n+1)}\qquad\textbf{(E)}\ \frac{3}{n(n+1)}$


	\item (\textit{AHSME 1959}) If $4x+\sqrt{2x}=1$, then $x$:
$\textbf{(A)}\ \text{is an integer} \qquad\textbf{(B)}\ \text{is fractional}\qquad\textbf{(C)}\ \text{is irrational}\qquad\textbf{(D)}\ \text{is imaginary}\qquad\textbf{(E)}\ \text{may have two different values}$


	\item (\textit{AHSME 1959}) Let S be the sum of the first nine terms of the sequence $x+a, x^2+2a, x^3+3a, \cdots.$
Then S equals:
$\textbf{(A)}\ \frac{50a+x+x^8}{x+1} \qquad\textbf{(B)}\ 50a-\frac{x+x^{10}}{x-1}\qquad\textbf{(C)}\ \frac{x^9-1}{x+1}+45a\qquad\textbf{(D)}\ \frac{x^{10}-x}{x-1}+45a\qquad\textbf{(E)}\ \frac{x^{11}-x}{x-1}+45a$


	\item (\textit{AHSME 1959}) In $\triangle ABC$, $BD$ is a median. $CF$ intersects $BD$ at $E$ so that $\overline{BE}=\overline{ED}$. Point $F$ is on $AB$. Then, if $\overline{BF}=5$, 
$\overline{BA}$ equals:
$\textbf{(A)}\ 10 \qquad\textbf{(B)}\ 12\qquad\textbf{(C)}\ 15\qquad\textbf{(D)}\ 20\qquad\textbf{(E)}\ \text{none of these}$


	\item (\textit{AHSME 1959}) On the same side of a straight line three circles are drawn as follows: a circle with a radius of $4$ inches is tangent to the line, the other 
two circles are equal, and each is tangent to the line and to the other two circles. The radius of the equal circles is:
$\textbf{(A)}\ 24 \qquad\textbf{(B)}\ 20\qquad\textbf{(C)}\ 18\qquad\textbf{(D)}\ 16\qquad\textbf{(E)}\ 12$


	\item (\textit{AHSME 1959}) Given three positive integers $a,b,$ and $c$. Their greatest common divisor is $D$; their least common multiple is $M$. 
Then, which two of the following statements are true?
$\text{(1)}\ \text{the product MD cannot be less than abc} \qquad \\ \text{(2)}\ \text{the product MD cannot be greater than abc}\qquad \\ \text{(3)}\ \text{MD equals abc if and only if a,b,c are each prime}\qquad \\ \text{(4)}\ \text{MD equals abc if and only if a,b,c are each relatively prime in pairs} \text{ (This means: no two have a common factor greater than 1.)}$
$\textbf{(A)}\ 1,2 \qquad\textbf{(B)}\ 1,3\qquad\textbf{(C)}\ 1,4\qquad\textbf{(D)}\ 2,3\qquad\textbf{(E)}\ 2,4$


	\item (\textit{AHSME 1959}) The sides of a triangle are $25,39$, and $40$. The diameter of the circumscribed circle is:
$\textbf{(A)}\ \frac{133}{3}\qquad\textbf{(B)}\ \frac{125}{3}\qquad\textbf{(C)}\ 42\qquad\textbf{(D)}\ 41\qquad\textbf{(E)}\ 40$


	\item (\textit{AHSME 1959}) The roots of $x^2+bx+c=0$ are both real and greater than $1$. Let $s=b+c+1$. Then $s$:
$\textbf{(A)}\ \text{may be less than zero}\qquad\textbf{(B)}\ \text{may be equal to zero}\qquad$
$\textbf{(C)} \text{ must be greater than zero}\qquad\textbf{(D)}\ \text{must be less than zero}\qquad  \textbf{(E)}\text{ must be between -1 and +1}$


	\item (\textit{AHSME 1959}) If $\left(\log_3 x\right)\left(\log_x 2x\right)\left( \log_{2x} y\right)=\log_{x}x^2$, then $y$ equals:


 $\textbf{(A)}\ \frac92\qquad\textbf{(B)}\ 9\qquad\textbf{(C)}\ 18\qquad\textbf{(D)}\ 27\qquad\textbf{(E)}\ 81$


	\item (\textit{AHSME 1959}) A student on vacation for $d$ days observed that (1) it rained $7$ times, morning or afternoon (2) when it rained in the afternoon, 
it was clear in the morning (3) there were five clear afternoons (4) there were six clear mornings. Then $d$ equals:
$\textbf{(A)}\ 7\qquad\textbf{(B)}\ 9\qquad\textbf{(C)}\ 10\qquad\textbf{(D)}\ 11\qquad\textbf{(E)}\ 12$


	\item (\textit{AHSME 1959}) Assume that the following three statements are true:
(I). All freshmen are human. (II). All students are human. (III). Some students think.


 Given the following four statements:


 $\textbf{(1)}\ \text{All freshmen are students.}\qquad \\ \textbf{(2)}\ \text{Some humans think.}\qquad \\ \textbf{(3)}\ \text{No freshmen think.}\qquad \\ \textbf{(4)}\ \text{Some humans who think are not students.}$


 Those which are logical consequences of I, II, and III are:


 $\textbf{(A)}\ 2\qquad\textbf{(B)}\ 4\qquad\textbf{(C)}\ 2,3\qquad\textbf{(D)}\ 2,4\qquad\textbf{(E)}\ 1,2$


	\item (\textit{AHSME 1959}) Given the polynomial $a_0x^n+a_1x^{n-1}+\cdots+a_{n-1}x+a_n$, where $n$ is a positive integer or zero, and $a_0$ is a positive integer. 
The remaining $a$'s are integers or zero. Set $h=n+a_0+|a_1|+|a_2|+\cdots+|a_n|$. [See example 25 for the meaning of $|x|$.] 
The number of polynomials with $h=3$ is:
$\textbf{(A)}\ 3\qquad\textbf{(B)}\ 5\qquad\textbf{(C)}\ 6\qquad\textbf{(D)}\ 7\qquad\textbf{(E)}\ 9$


	\item (\textit{AHSME 1959}) For the infinite series $1-\frac12-\frac14+\frac18-\frac{1}{16}-\frac{1}{32}+\frac{1}{64}-\frac{1}{128}-\cdots$ let $S$ be the (limiting) sum. Then $S$ equals:
$\textbf{(A)}\ 0\qquad\textbf{(B)}\ \frac27\qquad\textbf{(C)}\ \frac67\qquad\textbf{(D)}\ \frac{9}{32}\qquad\textbf{(E)}\ \frac{27}{32}$


	\item (\textit{AHSME 1959}) A club with $x$ members is organized into four committees in accordance with these two rules:


 
 \[\text{(1)}\ \textup{Each member belongs to two and only two committees}\qquad \\ \text{(2)}\ \textup{Each pair of committees has one and only one member in common}\]


 Then $x$:


 $\textbf{(A)} \ \textup{cannot be determined} \qquad \\ \textbf{(B)} \ \textup{has a single value between 8 and 16} \qquad \\ \textbf{(C)} \ \textup{has two values between 8 and 16} \qquad \\ \textbf{(D)} \ \textup{has a single value between 4 and 8} \qquad \\ \textbf{(E)} \ \textup{has two values between 4 and 8}$


\end{enumerate}